\scanpage{120}
% archive sheet number: 119
La donnée d'un isom $\kappa:\underline t_E\simeq\mathcal O_t(1)$ équivaut à celui de
\[
 \check\kappa:\underline{\check t}_E\simeq\underline{\mathcal O}_t(\unsure{+1|1})^\vee\simeq\underline{\mathcal O}_t(-1)
 \xrightarrow[\text{via }\omega]{}\underline{\mathcal O}_t(-1)(\underline\omega)\simeq F_t
\]
\deleted{ce qui implique une injection} \inserted{ou encore, à celui d'un} monomorphisme
\[
 \check t_E\hookrightarrow V_{\unsure{I|J}}(k)\hookrightarrow k^I
\]
dont l'image soit $F_t$, ou enfin à celui d'un élément non nul de
\[
 \deleted{\unsure{t_E}}\,V_I(t_E)=\deleted{\unsure{t_E^I}}\ \operatorname{Ker}(\underline t_E^I\xrightarrow[\text{somme}]{}\underline t_E)
\]
donc on peut considérer que la donnée de $\kappa$ revient à celle d'une famille
\[
 (\delta_i)_{i\in I}\in\underline t_E^I \tag*{(82)}
\]
satisfaisant certaines relations, savoir
\[
 \left\{\begin{aligned}
 &\text{a) }\sum\delta_i=0\quad\text{exprimant }\delta_i\in V_I(\underline t_E)\\
 &\text{b) Relation de proportionnalité, exprimant que}\\
 &\quad\text{l'image de }\check t_E\hookrightarrow V_I(k),\ \text{\unsure{défini|définie} par }(\delta_i)\text{ est }F_t,\\
 &\quad\text{de sorte que }(\delta_i)\text{ définit }\kappa:\underline t_E\simeq\check F_t\underset{\text{via }\omega}{\simeq}\mathcal O_t(1)\\
 &\text{c) Relation de l'épinglage de Legendre, exprimant}\\
 &\quad\text{que }\kappa^{\otimes2}\text{ est un isom. }\text{\unsure{standart|standard}}\text{ déjà connu.}
 \end{aligned}\right. \tag*{(83)}
\]
Nous allons expliciter les $\delta_i$ en termes du choix d'un $\delta\in\underline t_E^*$, \inserted{en plus \unsure{du|des} choix de $\omega\in\underline\omega$,} comme
\[
 \delta_i=\varphi_i^\omega(\delta)\delta\qquad\text{où}\qquad
 \begin{aligned}&\varphi_i^\omega(\delta)\in k^*\\ &(i\in I,\omega\in\underline\omega(I),\delta\in t_E^*)\end{aligned} \tag*{(84)}
\]
La condition d'invariance de $\delta_i$ p.r. au choix de $\delta$ s'exprime par la relation
\[
 \varphi_i^\omega(\lambda\delta)=\frac1\lambda\varphi_i^\omega(\delta) \tag*{(85)}
\]
